% ----------------------------------
%   Appendix C: Source code
% ----------------------------------

\chapter{Source code}
\label{app:code}

Every figure, table and number reported in this work is reproducible from the repository
\href{https://github.com/M315/TFM}{\texttt{github.com/M315/TFM}}.
This appendix gives the layout of the code and says which script produces what,
so that any result in the main text can be traced back to the routine that computed it.
Appendices \ref{app:impl_vol} and \ref{app:local_vol} discuss the choices made inside these
routines, here we only describe the structure.

The code is organized in two packages under \texttt{scripts/}.
The first builds the implied volatility surface from the raw option chain,
the second calibrates the local volatility surface from the prices that the first exports.
They run in separate conda environments,
because the finite element solver uses a legacy FEniCS stack without \texttt{pandas},
which is why the hand off between them is a file on disk rather than a function call.

\section{Implied volatility pipeline}
\label{subsec:code_impl_vol}

\begin{verbatim}
scripts/impl_vol/
    data.py           option chains from Yahoo Finance, cached by date
    rates.py          US Treasury par curve, bootstrapped to zero rates
    implied_vol.py    Black-Scholes prices, vega, implied vol inversion
    main.py           the implied volatility surface and smile slices
    diag_naive_q.py   the three stage dividend yield calibration
    export_smiles.py  exports the per expiration call price smiles
\end{verbatim}

The dependency flow is linear.
\texttt{data.py} and \texttt{rates.py} fetch and cache the raw inputs,
\texttt{implied\_vol.py} provides the pricing and inversion used by everything else,
and \texttt{main.py}, \texttt{diag\_naive\_q.py} and \texttt{export\_smiles.py} are the three
entry points.
The methodology behind them is Appendix \ref{app:impl_vol}.

\section{Local volatility calibration}
\label{subsec:code_local_vol}

\begin{verbatim}
scripts/local_vol_egger/
    pde/
        forward.py     forward Dupire solve, implicit Euler, P1 elements
        linearized.py  linearized (sensitivity) solve, one direction h
        adjoint.py     backward adjoint solve, accumulates the gradient
    optimization/
        vol.py         the coefficient a(y,tau) as a stack of spatial slices
        tikhonov.py    functional, gradient, and the L-BFGS-B driver
        surface.py     the time dependent extension, one slice per maturity
    utils.py           Black-Scholes helpers, meshes, smile handling
    derivative_test.py gradient checks (duality and Taylor)
    main.py            entry point, selects and runs one example
    examples/
        example1.py              Example 1, cases A to D
        example2.py              Example 2, oscillatory coefficient
        example3_market.py       single maturity fit on market data
        example4_surface.py      the calibrated market surface
        example5_skew.py         local vs implied skew, factor of two
        example6_holdout.py      the out of sample pricing test
        example7_arbitrage.py    density and calendar arbitrage check
        example8_discounting.py  the two normalizations of u
\end{verbatim}

The three modules in \texttt{pde/} are the three parabolic problems of
Section \ref{sec:gradient_adjoint}.
\texttt{optimization/tikhonov.py} assembles the functional \eqref{eq:tikhonov_functional} and
calls \texttt{pde/forward.py} and \texttt{pde/adjoint.py} once per iteration to obtain the
gradient \eqref{eq:tikhonov_gradient}.
\texttt{pde/linearized.py} is not used in the calibration itself,
it exists so that \texttt{derivative\_test.py} can check the adjoint against an independent
computation of the same directional derivative.
\texttt{optimization/surface.py} wraps the same machinery for the stack of slices
\eqref{eq:coefficient_slices} and minimizes \eqref{eq:surface_functional}.
Each file in \texttt{examples/} is self contained and writes its own figures.

\section{Where each result comes from}
\label{subsec:code_results_map}

\begin{table}[h!]
    \centering
    \begin{tabular}{l l}
        \toprule
        Result & Script \\
        \midrule
        Implied volatility surface and smiles, Chapter \ref{ch:bs}
            & \texttt{impl\_vol/main.py} \\
        Dividend yield calibration, Appendix \ref{app:impl_vol}
            & \texttt{impl\_vol/diag\_naive\_q.py} \\
        Call price smiles used as data
            & \texttt{impl\_vol/export\_smiles.py} \\
        Gradient verification, Table \ref{tab:taylor_test}
            & \texttt{local\_vol\_egger/derivative\_test.py} \\
        Example 1, Figures \ref{fig:egger_ex1_B} to \ref{fig:egger_ex1_BD}
            & \texttt{examples/example1.py} \\
        Example 2, Figure \ref{fig:egger_ex2}
            & \texttt{examples/example2.py} \\
        Market surface, Figures \ref{fig:egger_surface} to \ref{fig:egger_reprice}
            & \texttt{examples/example4\_surface.py} \\
        No-arbitrage check, Figure \ref{fig:egger_arbitrage}
            & \texttt{examples/example7\_arbitrage.py} \\
        Factor of two, Figure \ref{fig:egger_skew} and Table \ref{tab:egger_skew}
            & \texttt{examples/example5\_skew.py} \\
        Out of sample test, Figure \ref{fig:egger_holdout}
            & \texttt{examples/example6\_holdout.py} \\
        Discounted normalization, Appendix \ref{subsec:discounted_normalization}
            & \texttt{examples/example8\_discounting.py} \\
        \bottomrule
    \end{tabular}
    \caption{Each reported result and the script that produces it.}
    \label{tab:code_map}
\end{table}
